ALMA is a radio observatory with a revolutionary design. It consists of 66
high-precision antennas, operating in the wavelength range from 0.32\,mm to
8.60\,mm. The principal array has fifty antennas of 12\,m diameter each that
can work together as a single telescope in the so-called interferometric
mode. There is also another array of four 12\,m antennas, and twelve smaller
antennas of 7\,m diameter each.

Imagine that a single 12\,m antenna is being calibrated, pointing to a source
with a known incident flux of $1 \times 10^{-20}\,\mathrm{W/m}^2$.

\begin{description}
\item[(4.1)] Assuming that all the flux arrives at the shortest wavelength of
  ALMA sensitivity, determine the average number of photons that would reach
  the detector every second. \hfill[2.0\,pt]
\item[(4.2)] Compare it to the average number of photons that would have
  reached the detector, if all the flux arrived at the longest wavelength of
  operation. \hfill[2.0\,pt]
\item[(4.3)] What is the angular resolution (in arcsec) of a single 12\,m
  antenna, operating at 74.9\,GHz? \hfill[2.0\,pt]
\item[(4.4)] Imagine the principal array operating at 74.9\,GHz in the
  interferometric mode. Assuming for simplicity that the spatial resolution
  is solely given by the longest baseline (largest distance between any pair
  of antennas), which turns to be $D_{\max} = 16$\,km, what would be the
  angular resolution (in arcsec) in this case? Treat this case as a single
  slit aperture instead of a circular one. \hfill[2.0\,pt]
\item[(4.5)] For a radio antenna, the term SEFD refers to `System Equivalent
  Flux Density', which is a characteristic energy flux density of the
  antenna, depending on its temperature and size. We also note that for
  energy estimation of radio photons, Rayleigh-Jeans approximation is valid.
  Assuming a system temperature of 691\,K, what would be the SEFD of the full
  ALMA observatory in Jansky if all the 66 antennas could work together?
  \hfill[2.0\,pt]
\end{description}
